% Insert after Section G of the supplementary material.
% Table numbers follow the surrounding manuscript automatically.
% No additional LaTeX packages are required.
% Source: sources/component_p1.csv, original generation seeds 0/1/2.
% SHA-256: c2468fc11cb86d345fbc013b6acf3b6f36fed768a1755e6946e1c25f8bd1a52b

\subsection*{H. Component Ablation and Category-wise Sensitivity}
\label{sec:flash_component_ablation}

\paragraph{Experimental protocol}

We evaluate how selected FLASH components affect threshold calibration.
We use SuperADD with the DINOv3 ViT-H/16+ backbone and
$448\times448$ inputs on all eight MVTec AD 2 categories.
We compare six single-axis variants with the full-FLASH configuration.
Each variant changes one synthesis setting while retaining the original
three-source defect bank and the remaining generation parameters.
This study is separate from the three-category bank-size ablation.

The variants test foreground restriction, the placement field, adaptive
placement, color harmonization, and blending.
OBS off removes the foreground placement restriction but retains the
host-eligibility policy.
Color-only OBS retains the color cue and removes the texture cue.
The uniform-blob variant replaces MRSP with a compact, uniformly located
blob of matched object-relative area.
Original placement only sets the adaptive-placement fraction to zero,
instead of the default mixture of original and adaptive placement.
The remaining variants disable CIELAB harmonization or replace hybrid
blending with alpha blending for every sample.

Generation seeds 0, 1, and 2 are paired with detector seeds 1, 2, and 3.
Detectors are fitted using normal training images, with native validation
normals used for fit-time validation.
Within each category and seed, we evaluate the ablated configurations
using one shared fitted detector.
The full-FLASH reference is fitted separately with the same seed.
Baseline comparisons are therefore paired by category and seed, rather
than by an identical fitted detector instance.
Calibration does not update the detector representation.

For each configuration, we select separate image and pixel thresholds by
maximizing F1 on native validation normals and synthetic anomalies.
We keep the positive budget fixed across variants: 90 samples per category and
seed, except 60 for fruit jelly and 105 for vial.
These budgets match the public-test anomaly counts, so the budget choice
is test-informed.
The thresholds are transferred unchanged to the real public-test set.
Reported scores are not real-test-optimized F1 maxima.
We average seeds within each category and then weight the eight
categories equally.
The reported standard deviation is computed across the three seed-level
macro scores.
F1 scores use a 0--100 scale; all differences are absolute changes on this scale.

\paragraph{Aggregate component effects}

Table~\ref{tab:flash_component_macro} reports the aggregate results.
The re-evaluated full-FLASH reference obtains 38.61\% pixel F1 and
78.23\% image F1.
Note that these values are from a separate evaluation and differ from
the main-paper results.
We therefore compute all changes against this reference.

Removing the OBS restriction reduces pixel and image F1 by 1.71 and
4.62, respectively.
Replacing MRSP with a matched-area uniform blob reduces the same scores
by 0.74 and 3.04.
These comparisons support the full foreground restriction and structured
placement field at the aggregate level for this detector.
Using original placement alone reduces pixel F1 by 1.07, while image
F1 increases by 0.11.
The mixed placement policy therefore has a pixel-level advantage in this
comparison, but not a consistent advantage across both metrics.

The remaining variants reveal different trade-offs.
Disabling harmonization increases pixel F1 by 1.03.
Alpha-only blending increases it by 0.82.
Neither changes the mean image F1.
Color-only OBS gives the highest pixel F1 among the reported variants,
41.86\%, an increase of 3.24, with a 0.07 decrease in image F1.
Thus, the full configuration is not the best observed setting for every
calibration metric.

\begin{table*}[t]
    \centering
    \small
    \setlength{\tabcolsep}{6pt}
    \begin{tabular}{lrrrr}
        \hline
        \textbf{Configuration} & \textbf{Pixel F1 (\%)} &
        \textbf{$\Delta$ pixel F1} & \textbf{Image F1 (\%)} &
        \textbf{$\Delta$ image F1} \\
        \hline
        Full FLASH & $38.61 \pm 1.46$ & $+0.00$ & $78.23 \pm 0.11$ & $+0.00$ \\
        OBS off & $36.91 \pm 1.71$ & $-1.71$ & $73.61 \pm 1.83$ & $-4.62$ \\
        MRSP: uniform blob & $37.88 \pm 1.39$ & $-0.74$ & $75.19 \pm 1.79$ & $-3.04$ \\
        Original placement only & $37.54 \pm 0.78$ & $-1.07$ & $78.34 \pm 0.79$ & $+0.11$ \\
        Harmonization off & $39.64 \pm 0.76$ & $+1.03$ & $78.23 \pm 0.11$ & $+0.00$ \\
        Alpha-only blending & $39.44 \pm 1.16$ & $+0.82$ & $78.23 \pm 0.11$ & $+0.00$ \\
        Color-only OBS & $41.86 \pm 1.22$ & $+3.24$ & $78.15 \pm 0.15$ & $-0.07$ \\
        \hline
    \end{tabular}
    \caption{Single-axis component ablations on SuperADD.
    Scores are transferred-threshold real-test F1 (\%), macro-averaged over
    all eight MVTec AD 2 categories.
    The $\pm$ values are sample standard deviations across three seed-level
    macro scores, not confidence intervals.
    Differences use the current full-FLASH control and are computed before
    rounding.}
    \label{tab:flash_component_macro}
\end{table*}

\paragraph{Category-wise effects}

Table~\ref{tab:flash_component_category} presents category-wise pixel F1
changes for the same configurations.
This analysis reveals opposing effects that are obscured by the macro
average.
For example, removing OBS reduces pixel F1 by 24.88 on fruit jelly and
23.17 on vial.
It increases the score by 22.50 on sheet metal and 27.60 on
wallplugs.
These changes partly cancel, leaving a macro decrease of only 1.71.
OBS is therefore influential, but its effect is not uniformly positive
across categories.

The improvements from simpler settings are also category-dependent.
Color-only OBS gains 18.04 on sheet metal but loses 8.31 on
wallplugs and 2.06 on walnuts.
Alpha-only blending gains 8.26 on vial but loses 4.57 on walnuts.
These results do not identify a universally preferred cue or compositing
operator.
Calibration F1 also does not establish comparative visual quality.
The alpha-only comparison changes the compositing policy and therefore
cannot establish that the Poisson collar is unnecessary.

\begin{table*}[t]
    \centering
    \footnotesize
    \setlength{\tabcolsep}{4pt}
    \begin{tabular}{lrrrrrrrrr}
        \hline
        \textbf{Configuration} & \textbf{Can} & \textbf{Fabric} &
        \textbf{\shortstack{Fruit\\Jelly}} & \textbf{Rice} &
        \textbf{\shortstack{Sheet\\Metal}} & \textbf{Vial} &
        \textbf{Wallplugs} & \textbf{Walnuts} & \textbf{Macro} \\
        \hline
        Full FLASH (F1, \%) & $0.00$ & $69.12$ & $55.39$ & $52.55$ &
        $7.76$ & $38.91$ & $18.51$ & $66.67$ & $38.61$ \\
        \hline
        OBS off & $+0.00$ & $-16.44$ & $-24.88$ & $-2.78$ &
        $+22.50$ & $-23.17$ & $+27.60$ & $+3.51$ & $-1.71$ \\
        MRSP: uniform blob & $+0.00$ & $-13.73$ & $-0.39$ & $-7.26$ &
        $+5.98$ & $+2.60$ & $+5.60$ & $+1.29$ & $-0.74$ \\
        Original placement only & $+0.00$ & $+5.20$ & $+0.60$ & $-3.89$ &
        $-0.06$ & $-8.23$ & $-2.31$ & $+0.10$ & $-1.07$ \\
        Harmonization off & $+0.00$ & $-0.94$ & $+0.47$ & $+4.38$ &
        $-0.14$ & $+4.22$ & $-0.37$ & $+0.61$ & $+1.03$ \\
        Alpha-only blending & $+0.00$ & $+5.38$ & $-0.34$ & $-0.14$ &
        $-0.47$ & $+8.26$ & $-1.54$ & $-4.57$ & $+0.82$ \\
        Color-only OBS & $+0.00$ & $+8.83$ & $+0.40$ & $+5.88$ &
        $+18.04$ & $+3.17$ & $-8.31$ & $-2.06$ & $+3.24$ \\
        \hline
    \end{tabular}
    \caption{Category-wise pixel F1 effects for the configurations in
    Table~\ref{tab:flash_component_macro}.
    The first row gives absolute full-FLASH pixel F1 (\%).
    All subsequent rows give changes from that category's control,
    averaged over three seeds.
    Positive values indicate higher F1.
    The macro column weights all eight categories equally.}
    \label{tab:flash_component_category}
\end{table*}

\paragraph{Interpretation and limitations}

Together, the results indicate component sensitivity rather than a uniform benefit from
every stage.
Foreground restriction and structured placement improve aggregate F1 in
the reported comparisons, while harmonization and hybrid blending are
not required to obtain the highest observed pixel F1.
The category-wise reversals do not support a universally preferred
configuration.

Image F1 alone should not be interpreted as successful normal-image
rejection.
The control predicts every image as anomalous in 21 of the 24
category/seed cases.
A separate threshold audit of the same control conditions measures a
95.89\% macro normal-image false-positive rate.
Generation and detector seeds vary together, and baseline comparisons
include separate fits.
The observed spread therefore does not isolate synthesis uncertainty or
establish statistical significance.
We report a selected subset of the component study, and all categories
have been inspected.
The comparison is therefore descriptive rather than an independent
validation of a revised configuration.
It does not isolate component interactions or establish transfer to other
detectors.
